% LTeX: language=fr
% Copyright © 2014, Loïc Grobol <loic.grobol@gmail.com>
%
% Permission is granted to Do What The Fuck You Want To
% with this document.
%
% See the WTF Public License, Version 2 as published by Sam Hocevar
% at http://wtfpl.net if you need more details.

\RequirePackage{xparse}
% Settings
\newcommand\myname{L. Grobol}
\newcommand\mylab{}
\newcommand\coursetitle{Mathématiques pour le TAL 2}
\newcommand\examtitle{Devoir en temps libre}
\newcommand\pdftitle{\coursetitle : \examtitle}
\newcommand\mymail{lgrobol@parisnanterre.fr}
\newcommand\titlepagetitle{\pdftitle}
\newcommand\titlepagesubtitle{}
\newcommand\docdate{2026-03-26}  % chktex 8
\newcommand\conference{M1 Plurital}

%%%% Type de document %%%%
\documentclass[a4paper, 11pt]{exam}
\pagestyle{headandfoot}

%%%% Packages %%%%%
% Lang
\usepackage{polyglossia}
    \setmainlanguage{french}
	\setotherlanguage[variant=british]{english}

% Document and typesetting
\usepackage{fontspec}
\directlua{
	luaotfload.add_fallback(
		"myfallback",
		{
			"NotoColorEmoji:mode=harf;",
			"NotoSans:mode=harf;",
			"DejaVuSans:mode=harf;",
		}
	)
}
\setmainfont[
	UprightFont={* Regular},
	ItalicFont={* Italic},
	BoldFont={* SemiBold},
	BoldItalicFont={* SemiBold Italic},
	RawFeature={fallback=myfallback;multiscript=auto;},
]{Libertinus Serif}
\setsansfont[
	RawFeature={fallback=myfallback;multiscript=auto;}
	]{Libertinus Sans}
\setmonofont[
	Scale=MatchLowercase,
	RawFeature={fallback=myfallback;multiscript=auto;},
]{Fira Mono}
					
\usepackage{bookmark}
\usepackage{caption}
	\captionsetup{skip=1ex, labelformat=empty}
\usepackage[type={CC}, modifier={by}, version={4.0}]{doclicense}
\usepackage{emptypage}
\usepackage[inline]{enumitem}
% \usepackage  % Cool but messes with fleqn, activate if necessary{fnpct}  % Better display of footnote next to punctuation
\usepackage{microtype}
\usepackage{parskip}
% \usepackage[subtle]{savetrees}
\usepackage{siunitx}
\sisetup{
	detect-all,
	group-separator=\text{\,},
}
	\DeclareSIUnit{\quantity}{\relax}
	\DeclareSIUnit{\words}{mots}
	\DeclareSIUnit{\sentences}{phrases}
	% Needed for italics and bold numbers in siunitx S-aligned columns
	\robustify\itshape
	\robustify\bfseries
\usepackage{tabularray}

% Math
% Amsmath must be loaded before mathtools if it has options
\usepackage[fleqn]{amsmath}
% Mathtools should be loaded before unicode-math
\usepackage{mathtools}
	\newtagform{brackets}{[}{]}	% Pour des lignes d'équation numérotées entre crochets
	\mathtoolsset{showonlyrefs, showmanualtags, mathic}	% affiche les tags manuels (\tag et \tag*) et corrige le kerning des maths inline dans un bloc italique voir la doc de mathtools
	\usetagform{brackets}	% Utilise le style de tags défini plus haut

% Load this as early as possible in the maths section
\usepackage[
	math-style=french,
	warnings-off={
		mathtools-colon,
		mathtools-overbracket,
	},
]{unicode-math}
	\setmathfont{Libertinus Math}

\usepackage{amsthm}
\usepackage[d]{esvect}
\usepackage{interval}
\usepackage{lualatex-math}
\usepackage{mleftright}
\usepackage{newunicodechar}
	\newunicodechar{√}{\sqrt}

% \usepackage[backgroundcolor=white, bordercolor=orange]{todonotes}
%     \NewDocumentCommand\addlink{ O{} m }{\todo[color=green!40, #1]{#2}}

% \usepackage{tikz}
% 	% Colour palette from [Paul Tol's technical note](https://personal.sron.nl/~pault/data/colourschemes.pdf) v3.1
% 	% Bright scheme
% 	\definecolor{sronbrightblue}{RGB}{68, 119, 170}
% 	\definecolor{sronbrightcyan}{RGB}{102, 204, 238}
% 	\definecolor{sronbrightgreen}{RGB}{34, 136, 51}
% 	\definecolor{sronbrightyellow}{RGB}{204, 187, 68}
% 	\definecolor{sronbrightred}{RGB}{238, 102, 119}
% 	\definecolor{sronbrightpurple}{RGB}{170, 51, 119}
% 	\definecolor{sronbrightgrey}{RGB}{187, 187, 187}

% 	\definecolor{sronmutedindigo}{RGB}{51,34,136}

% 	% And my favourite purple
% 	\definecolor{myfavouritepurple}{RGB}{113, 10, 186}

% 	\NewDocumentCommand{\textnode}{O{}mm}{\tikz[remember picture, baseline=(#2.base), inner sep=0pt]{\node[#1] (#2) {#3};}}
% 	\NewDocumentCommand{\mathnode}{O{}mm}{\tikz[remember picture, baseline=(#2.base), inner sep=0pt]{\node[#1] (#2) {\(\displaystyle #3\)};}}
% 	% Global style
% 	\tikzset{
% 		>=stealth,
% 	}
% 	% Beamer utilities, for easy figure import
% 	\tikzset{
% 		visible on/.style={},
% 		accented/.style={text=accent, thick},
% 	}

% 	% Misc utilities
% 	\tikzset{
% 		true scale/.style={scale=#1, every node/.style={transform shape}},
% 	}

% 	\usetikzlibrary{tikzmark}
% 	\usetikzlibrary{matrix, chains}
% 	\usetikzlibrary{shapes, shapes.geometric}
% 	\usetikzlibrary{decorations.pathreplacing}
% 	\usetikzlibrary{positioning, calc, intersections}
% 	\usetikzlibrary{fit}

% 	% TikZ externalisation
% 	\usetikzlibrary{external}
% 	% Create the `tikzpics/` folder if it does not exist
% 	\usepackage{shellesc}
% 	\ShellEscape{mkdir _tikzpics}
% 	% Only externalise pictures on demand
% 	\tikzset{
% 		external/export=false,
% 		external/prefix=_tikzpics/,
% 	}
% 	\tikzexternalize

% \usepackage{forest}
% 	\useforestlibrary{linguistics}

% \usepackage{minted}
% 	\usemintedstyle{lovelace}
% 	\setminted{autogobble, tabsize=4}

\usepackage[
	english=american,
	french=guillemets,
	autostyle=true,
]{csquotes}
	\renewcommand{\mkbegdispquote}[2]{\itshape\let\emph\textbf}
	\renewcommand{\mkenddispquote}[2]{%
		\hfill\mbox{\upshape#1#2}%
	}
	% Like `\foreignquote` but use the outside language's quotes not the inside's
	\NewDocumentCommand\quoteforeign{m m}{\enquote{\textlang{#1}{#2}}}
	% \renewcommand{\mkbegdispquote}[2]{%
	%     \tikz[remember picture, overlay]{\fill[lightgray] (pic cs:quotestart) rectangle (pic cs:quoteend);}%
	%     \tikzmark{quotestart}%
	% }
	% \renewcommand{\mkenddispquote}[2]{%
	%     #1#2%
	%     \hfill\tikzmark{quoteend}%
	% }


% Annoying package loading order, so far we have:
% digraph packages {hyperref -> biblatex; hyperref -> hyperxmp; cleveref -> hyperref; fixme -> hyperref; bookmark -> hyperref}

% \usepackage[
% 	style=authoryear-comp,
% 	mergedate=basic,
% 	dashed=false,
% 	backend=biber,
% 	isbn=false,
% 	maxcitenames=2,
% 	uniquelist=false,
% 	maxbibnames=8,
% 	backref=true,
% ]{biblatex}
% 	% No small caps in french bib
% 	\DefineBibliographyExtras{french}{\restorecommand\mkbibnamefamily}
	% \addbibresource{biblio.bib}

\usepackage{hyperxmp}  % Load before hyperref
% \usepackage[a-3u]{pdfx}
% \usepackage[unicode]{hyperref}  % Apparently already loaded, no idea by whom. I hate you, latex
\hypersetup{
	pdftitle={\pdftitle},
	pdfauthor={\myname},
	pdfcontactemail={\mymail},
	colorlinks=true,
	breaklinks=true,
	hypertexnames=true,
	pdfdisplaydoctitle,
	pdflang={fr},
	keeppdfinfo,
}
% \usepackage{bookmark}  % Load after hyperref 
% \usepackage{footnotebackref}

% \usepackage[status=draft]{fixme} % Load after hyperref
% \usepackage[status=final]{fixme} % Load after hyperref
	% \fxsetup{theme=color}
\usepackage{cleveref}  % Load after hyperref


%%%% License %%%%


%%% Custom commands
% Maths
\NewDocumentCommand\transpose{m}{\prescript{\mathrm{t}}{}{#1}}

\DeclareMathOperator\Img{Im}
\DeclareMathOperator\Ker{Ker}
\DeclareMathOperator\FFNN{FFNN}
\DeclareMathOperator\LSTM{LSTM}
\NewDocumentCommand\bLSTM{}{\overleftarrow{\LSTM}}
\NewDocumentCommand\fLSTM{}{\overrightarrow{\LSTM}}
\DeclareMathOperator\GRU{GRU}
\NewDocumentCommand\bGRU{}{\overleftarrow{\GRU}}
\NewDocumentCommand\fGRU{}{\overrightarrow{\GRU}}
\DeclareMathOperator*\argmax{arg\,max}
\DeclareMathOperator*\softmax{softmax}

\DeclarePairedDelimiter\abs{\lvert}{\rvert}
\DeclarePairedDelimiter\brackets{[}{]}
\DeclarePairedDelimiterX\compset[2]{\lbrace}{\rbrace}{#1\,\delimsize|\,\mathopen{}#2}
\DeclarePairedDelimiterX\innprod[2]{\langle}{\rangle}{#1\,\delimsize|\,\mathopen{}#2}
\DeclarePairedDelimiter\norm{\lVert}{\rVert}
\DeclarePairedDelimiter\nnorm{\Vvert}{\Vvert}
\DeclarePairedDelimiter\discseg{⟦}{⟧}
\DeclarePairedDelimiter\ceil{⌈}{⌉}
\DeclarePairedDelimiter\floor{⌊}{⌋}
\DeclarePairedDelimiter\round{⌊}{⌉}

\NewDocumentCommand\fonct{ m m m m m }{
	\left|\begin{array}{rrl}
		#1: & #2 & ⟶ #3 \\
			& #4 & ⟼ #1\mleft(#4\mright) = #5
	\end{array}\right.
}

% Easy column vectors \vcord{a,b,c} ou \vcord[;]{a;b;c}
% Here be black magic
\ExplSyntaxOn
	\NewDocumentCommand{\vcord}{O{,}m}{\vector_main:nnnn{p}{\\}{#1}{#2}}
	\NewDocumentCommand{\tvcord}{O{,}m}{\vector_main:nnnn{psmall}{\\}{#1}{#2}}
	\seq_new:N\l__vector_arg_seq
	\cs_new_protected:Npn\vector_main:nnnn #1 #2 #3 #4{
		\seq_set_split:Nnn\l__vector_arg_seq{#3}{#4}
		\begin{#1matrix}
			\seq_use:Nnnn\l__vector_arg_seq{#2}{#2}{#2}
		\end{#1matrix}
	}

	\NewDocumentCommand{\concat}{sO{,}m}{
		\IfBooleanTF{#1}{
			\ensuremath{\brackets*{\concat_main:nnn{,~}{#2}{#3}}}
		}{
			\ensuremath{\brackets{\concat_main:nnn{,~}{#2}{#3}}}
		}
	}
	\seq_new:N\l__concat_arg_seq
	\cs_new_protected:Npn\concat_main:nnn #1 #2 #3 {
		\seq_set_split:Nnn\l__concat_arg_seq{#2}{#3}
		\seq_use:Nnnn\l__concat_arg_seq{#1}{#1}{#1}
	}
\ExplSyntaxOff
	
	
\let\card\abs

% Misc
% from https://tex.stackexchange.com/a/27099/8547
\makeatletter
\NewDocumentCommand\centerfloat{}{%
	\parindent \z@%
	\leftskip \z@ \@plus 1fil \@minus \textwidth%
	\rightskip\leftskip%
	\parfillskip \z@skip%
}
\makeatother

\NewDocumentCommand\definition{ o m }{\emph{#2}}

\title{\titlepagetitle}
\author{\myname}
\date{\docdate}

\firstpageheader{\bfseries\Large\coursetitle\\\examtitle}{}{\bfseries\Large\docdate}
% \runningheader{\coursetitle}{\examtitle}{\docdate}
\firstpagefooter{}{\thepage\ / \numpages}{}
\runningfooter{}{\thepage\ / \numpages}{}

%%%% Commandes spécifiques au document %%%%

\begin{document}
\pdfbookmark[1]{Cover}{cover}

\section{Valeur absolue}

\begin{definition}[Valeur absolue]
	On appelle \definition{valeur absolue de \(ℝ\)} la fonction
%
	\begin{equation}
		\left|\begin{array}{rrl}
			\abs{⋅}: & ℝ^n & ⟶ ℝ \\
				& x & ⟼ \abs{x}=
					\begin{dcases*}
				 		x & si \(x ⩾ 0\)\\
						-x & sinon
					\end{dcases*}
		\end{array}\right.
	\end{equation}
\end{definition}

Remarque : comme la définition de la valeur absolue est différente pour les réels positifs et les
réels négatifs, les démonstrations la concernant nécessitent souvent de distinguer plusieurs cas à
traiter indépendamment.

\begin{questions}
	\question Que valent \(\abs{2713}\), \(\abs{-1213}\), \(\abs{0}\) et \(\abs{19-91}\) ?
	\question Soit \(a∈ℝ\). En distinguant les cas \(a⩽0\) et \(a⩾0\), montrer que \(a⩽\abs{a}\) et \(\abs{-a} = \abs{a}\). En déduire que \(-a⩽\abs{a}\).
	\question Soit \(a∈ℝ\) et \(b∈ℝ\), montrer que \(\abs{ab} = \abs{a}\abs{b}\). On pourra distinguer le cas où \(a\) et \(b\) sont de même signe et le cas où ils sont de signes opposés.
	\question Soit \(a∈ℝ\) et \(b∈ℝ\). Montrer que\(\abs{a+b} ⩽ \abs{a} + \abs{b}\). On pourra distinguer les cas \(a+b ⩾ 0\) et \(a+b⩽0\).
	\question Soit \(a∈ℝ\) et \(b∈ℝ\) tels que \(b⩾a\) et \(b⩾-a\), montrer que \(b⩾\abs{a}\).
\end{questions}

\section{Normes vectorielles}

\begin{definition}[Norme vectorielle]
	On appelle \definition{norme vectorielle} sur un espace vectoriel \(E\) (par exemple \(ℝ^n\)), toute fonction \(N: E → ℝ\) telle que :

	\begin{enumerate}[label=(\roman*)]
		\item \(N\) est \emph{séparante} : pour tout \(u∈E\), \(N(u) ⩾ 0\), et \(N(u) = 0\) si et seulement si \(u=0\).
		\item \(N\) est \emph{absolument homogène} : pour tout \(u∈E\), pour tout \(λ∈ℝ\), \(N(λu) = \abs{λ}N(u)\).
		\item \(N\) est \emph{sous-additive} : pour tout \(u∈E\), pour tout \(v∈E\), on a \(N(u+v) ⩽ N(u) + N(v)\).
	\end{enumerate}
\end{definition}

\begin{questions}
	\question Montrer que la valeur absolue est une norme sur \(ℝ\).
	\question Soit \(N\), une norme sur \(E\), on veut montrer que la fonction suivante est une distance sur \(ℝ^n\) :
		%
		\begin{equation}
			\fonct{d_N}{E×E}{ℝ}{u, v}{\norm{v-u}}
		\end{equation}
		
		\begin{parts}
			\part Soit \(u∈E\), montrer que \(d_N(u, u) = 0\)
			\part Soit \(u∈E\) et \(v∈E\) tels que \(d_N(u,v) = 0\), montrer que \(u = v\). On rappel que deux vecteurs sont égaux si et seulement si \(u-v=0\).
			\part Soit \(u∈E\) et \(v∈E\), montrer que \(d_N(u,v)=d_N(v,u)\).
			\part Soit \(u∈E\), \(v∈E\) et \(w∈E\), montrer que \(d(u,w) ⩽ d(u,v) + d(v,w)\). On pourra utiliser le fait que \(u-w = (u-v) + (v-w)\). 
			\part Déduire des questions précédentes que \(d_N\) est une distance sur \(ℝ^n\).
		\end{parts}
	\question Soit \(u∈E\) et \(v∈E\)
		\begin{parts}
			\part Montrer que \(N(u-v) = N(v-u)\).
			\part En écrivant \(N(u-v) = N(u-v) + N(v) - N(v)\) et en utilisant la sous-additivité de \(N\), montrer  que \(N(u-v) ⩾ N(u) - N(v)\).
			\part Montrer que \(N(u-v) ⩾ N(v) - N(u)\) et en déduire que \(N(u-v) ⩾ \abs{N(u) - N(v)}\).
		\end{parts}
		\emph{Cette propriété est la \emph{seconde inégalité triangulaire}, elle peut aussi se montrer dans le cas plus général des distances.}
\end{questions}

\section{Norme euclidienne de \(ℝ^n\)}

\textbf{Rappel} : Pour tout \(n∈ℕ\), on appelle \definition{norme euclidienne de \(ℝ^n\)} la fonction
%
\begin{equation}
	\left|\begin{array}{rrl}
		\norm{⋅}_2: & ℝ^n & ⟶ ℝ \\
			& u & ⟼ \norm{u}_2 = √{\innprod{u}{u}}
	\end{array}\right.
\end{equation}

Quand il n'y a pas d'ambiguïté (comme dans cet exercice), on la note simplement \(\norm{u}\) (sans
le \(2\) en indice.

On veut montrer que la norme euclidienne de \(ℝ^n\) est bien une norme.

\begin{questions}
	\question Soit \(0_n\), le vecteur nul de \(ℝ^n\). Que vaut \(\norm{0_n}\) ?
	\question Soit \(v∈ℝ^n\), écrire \(\norm{v}²\) en fonction des coordonnées de \(v\). En déduire que \(\norm{v}\) est bien définie, que \(\norm{v}⩾0\) et que \(\norm{v}=0\) si et seulement si \(v=0\).
	\question
		\begin{parts}
			\part Démontrer que pour tout \(v∈ℝ^n\), \(\norm{-v}=\norm{v}\). On dit que \(\norm{⋅}\) est \emph{paire}.
			\part Démontrer que pour tout \(v∈ℝ^n\) et \(λ∈ℝ\), si \(λ⩾0\), alors, \(\norm{λv}=λ\norm{v}\). On dit que \(\norm{⋅}\) est \emph{positivement homogène}.
			\part Déduire des questions précédentes que pour tout \(v∈ℝ^n\) et \(λ∈ℝ\), \(\norm{λv}=\abs{λ}\norm{v}\).
		\end{parts}
	\question\footnote{Cette propriété est le théorème de Pythagore :)} Soit \(u∈ℝ^n\) et \(v∈ℝ^n\) tels que \(\innprod{u}{v}=0\). En développant \(\innprod{u+v}{u+v}\), montrer que \(\norm{u+v}²=\norm{u}²+\norm{v}²\).
	\question On veut montrer que pour tout \(u∈ℝ^n\) et \(v∈ℝ^n\), on a \(\innprod{u}{v}⩽\norm{u}\norm{v}\).
		\begin{parts}
			\part Démontrer l'inégalité ci-dessus dans le cas \(v=0\).
			\part Soit \(P\) la fonction définie par
				\begin{equation}
					\fonct{P}{ℝ}{ℝ}{t}{\norm{u+tv}²}
				\end{equation}
				Soit \(x∈ℝ\), que peut-on dire du signe de \(P(x)\) ?
			\part On suppose que \(v≠0\). Soit \(t₀=-\frac{\innprod{u}{v}}{\norm{v}²}\). Montrer en développant le produit scalaire que
				\begin{equation}
					P(t₀) = \norm{u}² - \frac{\innprod{u}{v}}{\norm{v}²}
				\end{equation}
			\part Déduire des questions précédentes que \(\innprod{u}{v}⩽\norm{u}\norm{v}\). Ce résultat s'appelle \emph{inégalité de Cauchy-Schwartz}. Rappel : si \(0 ⩽ x ⩽ y\), alors \(√{x} ⩽ √{y}\).
		\end{parts}
	\question En développant \(\innprod{u+v}{u+v}\), montrer que \(\norm{u+v}⩽\norm{u}+\norm{v}\).
	\question Déduire des questions précédentes que la norme euclidienne de \(ℝ^n\) est bien une norme sur \(ℝ^n\).
\end{questions}

\end{document}
